\documentclass[a4paper, UTF-8]{article}
\usepackage{algorithm}
\usepackage{algpseudocode}
\usepackage{amsmath}
\usepackage{amsfonts}
\usepackage{amsthm}
\usepackage{tikz}
\usetikzlibrary{automata, positioning, arrows, trees}
\usepackage[utf8]{inputenc}
\usepackage[T1]{fontenc}
\usepackage{hyperref}
\usepackage{url}
\usepackage{booktabs}
\usepackage{nicefrac}
\usepackage{microtype}
\usepackage{xcolor}
\usepackage{amssymb}

\theoremstyle{plain}
\newtheorem{theorem}{\indent Theorem}[section]
\newtheorem{lemma}[theorem]{\indent Lemma}
\newtheorem{assumption}{\indent Assumption}[section]
\newtheorem{note}[theorem]{\indent Notation}
\newtheorem{proposition}[theorem]{\indent Proposition}
\newtheorem{corollary}[theorem]{\indent Corollary}
\newtheorem{definition}{\indent Definition}[section]
\newtheorem{example}{\indent Example}[section]
\newtheorem{remark}{\indent Remark}[section]
\newenvironment{solution}{\begin{proof}[\indent\bf Solution]}{\end{proof}}
\renewcommand{\proofname}{\indent\bf Proof}

\begin{document}

\title{\bf{Lab 9 Answers}}
\author{Yanqiao Chen*\\ 
  Department of Computer Science and Engineering\\
  Southern University of Science and Technology\\
  \texttt{12412115@mail.sustech.edu.cn}
  }
\maketitle

\section{Question 1} 

\begin{itemize}
  \item No. For $n_1 = 1 \ne n_2 = 3$, we have $f(n_1) = f(n_3) = 6$. 
  \item No. For $m = 10$, there is no such $n$ such that $f(n) = m$. 
  \item No. From the previous two question we may know. 
  \item Yes. For any $n_1 \ne n_2$, $f(n_1) \ne f(n_2)$. 
  \item Yes. For any $m$, we have $n$ such that $f(n) = m$. 
  \item Yes.
\end{itemize}

\section{Question 2} 




\end{document}